\subsection{Inducing the marked $S_4$ quotient does not remove the harmonic frame ambiguity}

\paragraph{D23.F16.}
Combine the first-pair Krammer representation of D23.F13
with the induced $24$-point representation of
$\widehat B=C_2\amalg C_3$ from D23.C11, and retain
the signed projective level-twelve quotient as well.
This combination is formed as a representation of
$\widehat B_3$, pulling the latter two factors back
along $\widehat B_3\to\widehat B$. In its Krammer
coordinate, ``centralizing $s$'' means centralizing
the half-twist lift $s_P=K_1K_2K_1$, whose square
is central; the six-dimensional representation is
not asserted to factor through $\widehat B$.
The finite centralizer of $s$ in the image of $F'$
still contains an element $\beta$ with
\[
 q_4(\beta)=I,\qquad q_3(\beta)=S
 =\begin{pmatrix}0&1\\-1&0\end{pmatrix}.
\tag{D23.F16.1}
\]
Its induced $24$-point coordinate can be the identity.
Thus this combined quotient does not determine the
harmonic parity from the finite harmonic frame equation,
even after imposing $\beta\in\operatorname{im}(F')$.
This is a finite ambiguity, not an assertion that the
nonidentity finite conjugator is an actual pure frame.

\begin{proof}
Let $P$ be the Krammer first-pair group.  The exact
derived orders of D23.P6 are
$46656,3888,486,243,3,1$. Thus $P'$ is solvable
and its abelianization has order $8$.
Let $E$ be the restriction to $F$ of the induced
representation in D23.C11. Its exact derived orders
are $6912,1728,64,1$, so $E'$ is solvable and its
abelianization has order $27$.
The simultaneous image of $F'$ in $P'\times E'$
is subdirect. A nontrivial common quotient of these
solvable groups would have a nontrivial abelianization,
which would be a quotient of both a two-group and
a three-group. This is impossible. Goursat's lemma
therefore gives the full product $P'\times E'$.

Under the diagonal-block projection of D23.F13,
\[
 P'\longrightarrow E_F'\times D_8'
\]
is surjective, with central kernel $\langle n\rangle$
of order three.  In the affine coordinates of D23.P3
choose the first coordinate
$\gamma=(0,S)\in k^2\rtimes Q_8=E_F'$,
and choose the second coordinate to be identity.
It follows that there is $\beta_P\in P'$ with
these diagonal values.

Put $s_P=K_1K_2K_1$. Its affine $W$ coordinate is
$(0,iS)$, which commutes with $\gamma$; its quotient
coordinate commutes with identity. Thus
$[s_P,\beta_P]\in\langle n\rangle$.
This kernel is central in $\langle K_1,K_2\rangle$,
and $s_P^2=(K_1K_2)^3$ is also central there.
Writing $s_P\beta_Ps_P^{-1}=\beta_Pn^j$ and
conjugating once more gives $n^{2j}=1$.
Since $n$ has order three, $j=0$. Every such lift
already centralizes $s_P$.

By the full-product statement, the pair
$(\beta_P,1)$ is the image of an element of $F'$.
It centralizes $s$ in both representations.  The
signed level-twelve coordinate of an element of
$F'$ has identity component modulo four: the two
signed generators $U^2,L^2$ commute modulo four,
so their commutator subgroup is killed there.
On $F'$ the linear part of the affine detector is
exactly $q_3$, since the modulus-four character
occurring in D23.P3 is trivial. Our choice of
$\gamma$ therefore gives $q_3(\beta_P)=S$.
This also shows that adjoining the level-twelve
coordinate imposes no further restriction on the
chosen pair: on $F'$ it is already determined by
the Krammer linear coordinate.  The resulting
matrix is
\[
 q_{12}(\beta)=\begin{pmatrix}9&4\\8&9\end{pmatrix},
\]
whose two reductions are those of (D23.F16.1).

For a literal finite certificate, the program
\texttt{krammer\_induced\_frame\_ambiguity.py}
finds the three lifts of $(\gamma,1)$ inside the
enumerated $P'$, checks that all commute with
$s_P$, and chooses the unique lift of order four.
Its complete six-dimensional matrix is in the
companion JSON. The existence proof above does
not depend on a guessed lift or on an assumed
congruence-kernel invariance.

At the identity action the harmonic frame equation
is simply that the frame conjugator centralizes $s$.
Both the identity and the constructed finite
conjugator solve this equation and satisfy the
positive modulus-four normalization. The latter
has odd harmonic residue modulo three, whereas
the former has even residue. Thus the finite
equation does not recover the actual parity.
The actual pure centralizer in the profinite
source is trivial; our assertion concerns only
the specified finite quotient.
\end{proof}
